\chapter{Filtracje i obserwacja procesów w czasie ciągłym}
\label{ch:filtracje}

\section{Informacja: co wolno wiedzieć w chwili \texorpdfstring{$t$}{t}?}
Proces $X=(X_t)_{t\ge0}$ opisuje, co dzieje się z układem.
Do zdefiniowania decyzji zależnej od przeszłości potrzebujemy jednak
osobnego pojęcia \emph{informacji dostępnej w danej chwili}.
Obserwator widzący ruch cząstki do czasu $t$ może obliczyć jej
aktualne położenie, maksymalną dotychczasową odległość lub moment
pierwszego przekroczenia progu. Nie powinien natomiast wiedzieć,
czy cząstka przekroczy próg \emph{jutro}.

Matematycznym modelem tego rozróżnienia jest rosnąca rodzina
$\sigma$-ciał: im późniejsza chwila, tym więcej zdarzeń umiemy
rozpoznać na podstawie dostępnej informacji.

\begin{definicja}[Filtracja i przestrzeń filtrowana]
Niech $(\Omega,\calF,\Prob)$ będzie przestrzenią
probabilistyczną, a $T=[0,\infty)$ lub $T=[0,T_0]$
przy $T_0>0$. \emph{Filtracją} nazywamy rodzinę
$\mathbb F=(\calF_t)_{t\in T}$ pod-$\sigma$-ciał $\calF$
spełniającą $\calF_s\subseteq\calF_t$ dla wszystkich
$s,t\in T$ z $s\le t$. Układ
$(\Omega,\calF,\mathbb F,\Prob)$ nazywamy
\emph{przestrzenią probabilistyczną z filtracją}.
\end{definicja}

\begin{definicja}[Filtracja naturalna i adaptacja]
Niech $X=(X_t)_{t\in T}$ będzie procesem o stanach w
$(E,\calE)$. Jego \emph{filtracją naturalną} jest
$\calF_t^X:=\sigma(X_s:0\le s\le t)$.
Mówimy, że $X$ jest \emph{adaptowany} do filtracji
$\mathbb F$, jeżeli $X_t$ jest $\calF_t/\calE$-mierzalny
dla każdego $t\in T$.
\end{definicja}

\begin{stwierdzenie}
Dla każdego procesu $X$ z czasem $T=[0,\infty)$
jego filtracja naturalna $\mathbb F^X$ jest filtracją,
$X$ jest do niej adaptowany, a jeżeli $\mathbb F$
jest dowolną filtracją, do której $X$ jest adaptowany,
to $\calF_t^X\subseteq\calF_t$ dla każdego $t\ge0$.
\end{stwierdzenie}
\begin{proof}
Jeżeli $s\le t$, to rodzina zmiennych
$\{X_u:0\le u\le s\}$ jest podrodziną
$\{X_u:0\le u\le t\}$, a więc
$\calF_s^X\subseteq\calF_t^X$.
Z definicji $X_t$ jest mierzalna względem
$\calF_t^X$. Dla ostatniej własności
adaptacja do $\mathbb F$ oznacza, że wszystkie
$X_u$, $u\le t$, są $\calF_u$-mierzalne,
a przez monotoniczność filtracji również
$\calF_t$-mierzalne. Zatem generowane przez nie
$\sigma$-ciało zawiera się w $\calF_t$.
\end{proof}

\begin{przyklad}[Losowe obserwacje w chwilach całkowitych]
Niech $\xi_1,\xi_2,\ldots$ będą niezależnymi rzutami symetryczną
monetą, o wartościach $\pm1$, i niech
$S_n=\sum_{k=1}^n\xi_k$, $S_0=0$.
Możemy rozważać $X_t=S_{\lfloor t\rfloor}$ dla $t\ge0$.
Filtracja $\calF_t:=\sigma(\xi_1,\ldots,\xi_{\lfloor t\rfloor})$
(z trywialnym $\calF_t$ dla $t<1$) opisuje poznawanie
wyników kolejnych rzutów. Dla $1\le t<2$ widzimy
pierwszy rzut, ale nie drugi. Proces $X$ jest adaptowany.
Zmienna $\xi_2$ nie jest $\calF_1$-mierzalna:
gdyby była, zdarzenie $\{\xi_2=1\}$ należałoby do
$\sigma(\xi_1)$, choć jest od niej niezależne i ma
prawdopodobieństwo $1/2$. Zdarzenie niezależne od
samego siebie miałoby prawdopodobieństwo $0$ lub $1$,
co daje sprzeczność.
\end{przyklad}

\begin{center}
\begin{tikzpicture}[>=Latex]
\draw[->] (0,0)--(12,0) node[right] {$t$};
\foreach \x/\lab in {1/0,4/1,7/2,10/3}{
  \draw (\x,0.12)--(\x,-0.12) node[below] {$\lab$};
}
\draw[very thick,blue!70!black] (1,1.1)--(4,1.1);
\draw[very thick,blue!70!black] (4,1.55)--(7,1.55);
\draw[very thick,blue!70!black] (7,0.68)--(10,0.68);
\draw[blue!70!black,fill=white,thick] (4,1.1) circle (2pt) (7,1.55) circle (2pt);
\fill[blue!70!black] (1,1.1) circle (2pt) (4,1.55) circle (2pt) (7,0.68) circle (2pt);
\node[left,blue!65!black] at (1,1.12) {$X_t$};
\node[draw,rounded corners,fill=green!8,align=center] at (2.5,2.45) {$\calF_0$\\bez rzutów};
\node[draw,rounded corners,fill=green!8,align=center] at (5.5,2.45) {$\calF_1$\\po I rzucie};
\node[draw,rounded corners,fill=green!8,align=center] at (8.5,2.45) {$\calF_2$\\po II rzucie};
\draw[->,green!45!black] (3.55,2.45)--(4.4,2.45);
\draw[->,green!45!black] (6.55,2.45)--(7.4,2.45);
\node[below=18pt,align=center] at (5.9,-0.2) {informacja rośnie, a przyszły przyrost\\nie jest znany przed wykonaniem rzutu};
\end{tikzpicture}
\par\smallskip\small Rozróżnienie: wykres trajektorii oraz zbiory informacji $\calF_t$.
\end{center}

\section{Zupełność i prawostronna ciągłość filtracji}
\begin{definicja}[Warunki zwykłe]
Niech $\mathbb F=(\calF_t)_{t\ge0}$ będzie filtracją
na $(\Omega,\calF,\Prob)$. Mówimy, że filtracja jest
\emph{zupełna}, jeśli $\calF_0$ zawiera wszystkie
podzbiory zbiorów $N\in\calF$ z $\Prob(N)=0$
(przyjmujemy, że przestrzeń bazowa jest odpowiednio uzupełniona).
Mówimy, że jest \emph{prawostronnie ciągła}, jeśli
dla każdego $t\ge0$ zachodzi
\[
\calF_t=\calF_{t+}:=\bigcap_{u>t}\calF_u.
\]
Filtracja spełnia \emph{warunki zwykłe}, jeśli jest
zupełna i prawostronnie ciągła.
\end{definicja}

Skąd prawostronna ciągłość? Jeśli informacja, że pewne
zdarzenie zaszło, jest dostępna w każdej chwili
$t+\varepsilon$ dla dowolnie małego $\varepsilon>0$,
to naturalnie chcielibyśmy traktować ją jako dostępną
już w chwili $t$. To założenie ma realne konsekwencje:
pomaga wykazać, że czasy pierwszego wejścia do
zbiorów otwartych są czasami zatrzymania.
Nie jest jednak automatyczną własnością każdej filtracji.

\begin{stwierdzenie}[Prawostronnie ciągłe powiększenie]
Niech $\mathbb F^0=(\calF_t^0)_{t\ge0}$ będzie
filtracją na pełnej przestrzeni probabilistycznej
$(\Omega,\calF,\Prob)$, a $\calN$ rodziną podzbiorów
zdarzeń zerowych z $\calF$. Niech
$\calH_t:=\sigma(\calF_t^0\cup\calN)$ i
$\calG_t:=\bigcap_{u>t}\calH_u$.
Wówczas $(\calG_t)_{t\ge0}$ jest zupełną,
prawostronnie ciągłą filtracją oraz
$\calF_t^0\subseteq\calG_t$ dla każdego $t$.
\end{stwierdzenie}
\begin{proof}
Rodzina $\calH_t$ jest rosnąca, dlatego również
$\calG_t$ jest rosnąca. Każde $\calG_t$ zawiera
$\calH_t$, ponieważ $\calH_t\subseteq\calH_u$
dla $u>t$, a zatem zawiera $\calF_t^0$ i $\calN$.
Dla $t\ge0$ obliczamy
\[
\bigcap_{s>t}\calG_s
=\bigcap_{s>t}\bigcap_{u>s}\calH_u
=\bigcap_{u>t}\calH_u=\calG_t.
\]
Środkowa równość wynika stąd, że dla każdego $u>t$
można wybrać $s$ z $t<s<u$, a warunki indeksowane
przez wszystkie takie pary $(s,u)$ są więc dokładnie
warunkami indeksowanymi przez $u>t$.
To dowodzi prawostronnej ciągłości i zupełności.
\end{proof}

\begin{uwaga}
Uzupełnienie i prawostronne domknięcie filtracji
nie zmienia samego procesu $X_t$ jako rodziny funkcji.
Może jednak zmienić, jakie zdarzenia uważa się za
obserwowalne w chwili $t$. Dlatego przy twierdzeniach
odwołujących się do warunków zwykłych trzeba je
wymienić jawnie, a nie zakładać milcząco.
\end{uwaga}


\begin{przyklad}[Filtracja naturalna nie musi być prawostronnie ciągła]
Niech $Z$ będzie niedegenerowaną zmienną Bernoulliego,
a $X_0=0$, $X_t=Z$ dla każdego $t>0$.
Wówczas $\calF_0^X=\{\varnothing,\Omega\}$, lecz
$\calF_t^X=\sigma(Z)$ dla każdego $t>0$, więc
\[
\calF_{0+}^X=\bigcap_{u>0}\calF_u^X
=\sigma(Z)\ne\calF_0^X.
\]
Wynik losowania $Z$ zostaje ujawniony natychmiast
\emph{po} chwili zerowej, choć w chwili $0$ nie jest
jeszcze obserwowalny. Prawostronne domknięcie filtracji
przenosi tę informację do chwili $0$.
Nie twierdzimy tutaj, że $X$ ma prawostronnie ciągłe
trajektorie w zerze: właśnie skok przy $t=0$ umożliwia
tak prosty kontrprzykład.
\end{przyklad}

\section{Mierzalność wspólna, progresywna i przewidywalna}
\begin{definicja}[Mierzalność procesu]
Niech $E$ będzie przestrzenią polską,
$X:[0,\infty)\times\Omega\to E$ oraz
$\mathbb F=(\calF_t)_{t\ge0}$ filtracją.
\begin{enumerate}
\item $X$ jest \emph{wspólnie mierzalny}, jeśli
$(t,\omega)\mapsto X_t(\omega)$ jest mierzalny względem
$\calB([0,\infty))\otimes\calF$.
\item $X$ jest \emph{progresywnie mierzalny}, jeśli dla każdego
$T>0$ jego ograniczenie do $[0,T]\times\Omega$
jest mierzalne względem
$\calB([0,T])\otimes\calF_T$.
\item \emph{Przewidywalne $\sigma$-ciało} $\calP$ na
$[0,\infty)\times\Omega$ jest $\sigma$-ciałem generowanym
przez wszystkie rzeczywiste procesy adaptowane,
których trajektorie są lewostronnie ciągłe na $(0,\infty)$
(ich wartości w $0$ są $\calF_0$-mierzalne).
Proces $E$-wartościowy nazywamy \emph{przewidywalnym},
jeśli jest mierzalny względem $\calP$ i $\calB(E)$.
\end{enumerate}
\end{definicja}

Znaczenie przewidywalności jest operacyjne: decyzja
realizowana w krótkim przedziale \emph{po chwili $t$}
powinna być wyznaczona informacją dostępną \emph{do}
chwili $t$, a nie skokiem, który nastąpi dopiero w tym
przedziale. Dlatego procesy proste w rachunku Itô
będą zbudowane z wartości znanych na lewych końcach
przedziałów.

\begin{twierdzenie}[Regularne trajektorie dają mierzalność progresywną]
\label{thm:prog}
Niech $E$ będzie przestrzenią polską,
$\mathbb F=(\calF_t)_{t\ge0}$ filtracją, a
$X=(X_t)_{t\ge0}$ procesem adaptowanym.
Jeśli każda trajektoria $t\mapsto X_t(\omega)$
jest prawostronnie ciągła, to $X$ jest progresywnie
mierzalny. To samo zachodzi przy lewostronnej
ciągłości na $(0,\infty)$ oraz dowolnej wartości w
$t=0$. W szczególności proces adaptowany o ciągłych
trajektoriach jest progresywnie mierzalny i przewidywalny.
\end{twierdzenie}
\begin{proof}
Ustalmy $T>0$. Dla każdej liczby naturalnej $n$
rozpatrzmy podział $[0,T]$ o węzłach
$t_k=kT/n$, $k=0,\ldots,n$. Jeśli trajektorie są
prawostronnie ciągłe, zdefiniujmy funkcję schodkową
\[
X^{(n)}_s(\omega)=
\begin{cases}
X_{t_k}(\omega),&t_{k-1}<s\le t_k,\quad 1\le k\le n,\\
X_0(\omega),&s=0.
\end{cases}
\]
Każda funkcja $X^{(n)}$ jest mierzalna względem
$\calB([0,T])\otimes\calF_T$, ponieważ jest sumą
skończenie wielu funkcji o stałej po czasie wartości
$X_{t_k}$ (w przypadku ogólnego $E$ oznacza to
funkcję sklejoną na skończonym podziale mierzalnym).
Dla każdego $s<T$ prawy koniec odpowiedniej komórki
zbiega do $s$ od strony prawej, toteż przez
prawostronną ciągłość $X^{(n)}_s\to X_s$;
dla $s=T$ równość zachodzi dokładnie, a dla $s=0$
równość jest w definicji. Granica punktowa funkcji
mierzalnych o wartościach w przestrzeni metrycznej
jest mierzalna, więc $X$ jest progresywnie mierzalny.

Jeżeli trajektorie są lewostronnie ciągłe na $(0,T]$,
zamiast prawych końców używamy lewych końców
komórek: $X^{(n)}_s=X_{t_{k-1}}$ dla
$t_{k-1}\le s<t_k$, a w chwili $T$ ustawiamy
$X^{(n)}_T=X_T$. Dla $s>0$ czasy aproksymacji
zbiegają do $s$ z lewej, a w chwili $0$ mamy
$X^{(n)}_0=X_0$. To dowodzi drugiej wersji.
Proces ciągły adaptowany jest z definicji
przewidywalny, ponieważ jest lewostronnie ciągły;
pierwsza część dowodu daje jego progresywną mierzalność.
\end{proof}

\begin{uwaga}[Dlaczego lewa aproksymacja wymaga ostrożności?]
W powyższym dowodzie lewostronnej ciągłości lewy koniec
komórki dla $s\in(t_{k-1},t_k)$ zbiega do $s$ z lewej.
Dla samego punktu siatki $s=t_k$ wartość
aproksymacji może być równa $X_s$, ale nie ma to
znaczenia dla granicy. Procesy jedynie adaptowane,
bez regularności trajektorii, nie muszą być wspólnie,
a tym bardziej progresywnie mierzalne.
\end{uwaga}

\begin{stwierdzenie}[Proste strategie przewidywalne]
Niech $0=t_0<t_1<\cdots<t_m<\infty$, a
$H_0,\ldots,H_{m-1}$ będą rzeczywistymi zmiennymi
losowymi, przy czym $H_k$ jest $\calF_{t_k}$-mierzalna.
Wtedy proces
\[
H_t=H_{-1}\ind{\{0\}}(t)
+\sum_{k=0}^{m-1}H_k\ind{(t_k,t_{k+1}]}(t),
\qquad 0\le t\le t_m,
\]
gdzie $H_{-1}$ jest dowolną $\calF_0$-mierzalną
zmienną losową, jest przewidywalny.
\end{stwierdzenie}
\begin{proof}
Dla $t\in(t_k,t_{k+1}]$ wartość $H_t=H_k$ jest
$\calF_{t_k}\subseteq\calF_t$-mierzalna.
W punkcie $t_k$ ($k\ge1$) wartość z poprzedniego
przedziału wynosi $H_{k-1}$ i też jest mierzalna
względem $\calF_{t_k}$. Proces jest więc adaptowany.
Przy każdym $t>0$ jego wartość jest granicą wartości
z czasów $s\uparrow t$, ponieważ przedziały mają
postać $(t_k,t_{k+1}]$. Trajektorie są lewostronnie
ciągłe, a więc $H$ jest przewidywalny z definicji.
\end{proof}

\begin{uwaga}
Zapis $H_{-1}$ nie oznacza informacji w chwili ujemnej;
służy wyłącznie do oddzielnego określenia wartości w $0$.
Przy konstrukcji całki Itô wybierzemy zwykle $H_0=0$
i pominiemy wkład pojedynczej chwili $0$.
\end{uwaga}

\section{Czasy zatrzymania}
\begin{definicja}[Czas zatrzymania]
Niech $\mathbb F=(\calF_t)_{t\ge0}$ będzie filtracją.
Zmienną losową $\tau:\Omega\to[0,\infty]$
nazywamy \emph{czasem zatrzymania względem $\mathbb F$},
jeżeli dla każdego $t\ge0$ zachodzi
$\{\tau\le t\}\in\calF_t$.
Jeśli dodatkowo $\Prob(\tau<\infty)=1$, to
mówimy, że $\tau$ jest prawie na pewno skończony.
\end{definicja}

Warunek mówi, że w chwili $t$ umiemy rozstrzygnąć,
czy zatrzymanie nastąpiło \emph{najpóźniej do $t$}.
Wymóg ten wyklucza strategie, które wybierają czas
na podstawie wyniku z przyszłości. Nie wymaga natomiast,
aby w chwili $0$ było wiadomo, jaka będzie wartość $\tau$.

\begin{stwierdzenie}[Działania na czasach zatrzymania]
Niech $\sigma,\tau$ będą czasami zatrzymania dla filtracji
$\mathbb F=(\calF_t)_{t\ge0}$. Wtedy
$\sigma\wedge\tau:=\min(\sigma,\tau)$,
$\sigma\vee\tau:=\max(\sigma,\tau)$ oraz
$\tau\wedge a$ dla każdej stałej $a\ge0$ są czasami
zatrzymania. Jeśli $\mathbb F$ jest prawostronnie ciągła,
to również $\inf_{n\ge1}\tau_n$ jest czasem zatrzymania
dla dowolnego ciągu czasów zatrzymania $(\tau_n)_{n\ge1}$.
\end{stwierdzenie}
\begin{proof}
Dla wszystkich $t\ge0$ mamy
$\{\sigma\wedge\tau\le t\}
=\{\sigma\le t\}\cup\{\tau\le t\}$ oraz
$\{\sigma\vee\tau\le t\}
=\{\sigma\le t\}\cap\{\tau\le t\}$.
Oba zdarzenia są w $\calF_t$.
Dla $\tau\wedge a$ stosujemy pierwszy przypadek,
bo stała $a$ jest czasem zatrzymania.

Niech $\rho=\inf_n\tau_n$. Dla dowolnego $u>0$
$\{\rho<u\}=\bigcup_{n\ge1}\{\tau_n<u\}$:
jeśli infimum jest mniejsze od $u$, to pewien
składnik jest mniejszy od $u$.
Każde $\{\tau_n<u\}$ należy do $\calF_u$,
bo jest przeliczalną sumą zbiorów
$\{\tau_n\le q\}$ po wymiernych $q<u$.
Stąd $\{\rho<u\}\in\calF_u$.
Ponieważ
$\{\rho\le t\}=\bigcap_{k\ge1}\{\rho<t+1/k\}$,
to zdarzenie po lewej należy do
$\bigcap_{k\ge1}\calF_{t+1/k}=\calF_t$
przy prawostronnej ciągłości filtracji.
\end{proof}

\begin{twierdzenie}[Czas pierwszego trafienia w zbiór domknięty]
\label{thm:hitting-closed}
Niech $(E,d)$ będzie przestrzenią metryczną,
$F\subseteq E$ niepustym zbiorem domkniętym,
$\mathbb F=(\calF_t)_{t\ge0}$ filtracją, a
$X=(X_t)_{t\ge0}$ procesem $E$-wartościowym,
adaptowanym, którego wszystkie trajektorie są ciągłe.
Definiujemy $\tau_F(\omega)=\inf\{t\ge0:X_t(\omega)\in F\}$,
przy $\inf\varnothing=\infty$.
Wówczas $\tau_F$ jest czasem zatrzymania.
\end{twierdzenie}
\begin{proof}
Dla każdego $t\ge0$ funkcja
$u\mapsto d(X_u(\omega),F)$ jest ciągła na
przedziale zwartym $[0,t]$; odległość od zbioru
domkniętego jest funkcją $1$-Lipschitzowską.
Wobec tego
\[
\{\tau_F\le t\}
=\left\{\inf_{0\le u\le t}d(X_u,F)=0\right\}
=\left\{\inf_{u\in D_t}d(X_u,F)=0\right\},
\]
gdzie $D_t=(\Q\cap[0,t])\cup\{0,t\}$ jest przeliczalny
i gęsty w $[0,t]$ (dla $t=0$ zawiera punkt $0$).
Pierwsza równość wymaga uzasadnienia w obie strony.
Jeżeli proces trafia do $F$ przed $t$ lub w chwili $t$,
infimum odległości wynosi $0$. Jeśli infimum wynosi $0$,
to na mocy zwartości $[0,t]$ jest osiągane w pewnej
chwili $u\in[0,t]$; domkniętość $F$ implikuje
$X_u\in F$, więc $\tau_F\le t$.
Druga równość wynika z ciągłości funkcji odległości
i gęstości $D_t$.
Wszystkie $X_u$ dla $u\in D_t$ są
$\calF_u\subseteq\calF_t$-mierzalne,
więc infimum przeliczalnej rodziny ich odległości
jest $\calF_t$-mierzalne. Zatem
$\{\tau_F\le t\}\in\calF_t$.
\end{proof}

\begin{twierdzenie}[Czas wejścia do zbioru otwartego]
Niech $(E,d)$ będzie przestrzenią metryczną,
$O\subseteq E$ otwartym zbiorem,
$\mathbb F=(\calF_t)_{t\ge0}$ prawostronnie ciągłą filtracją,
a $X=(X_t)_{t\ge0}$ adaptowanym procesem o ciągłych
trajektoriach. Wtedy
$\tau_O:=\inf\{t\ge0:X_t\in O\}$ (z $\inf\varnothing=\infty$)
jest czasem zatrzymania.
\end{twierdzenie}
\begin{proof}
Dla każdego $t>0$ mamy
$\{\tau_O<t\}=\bigcup_{q\in\Q\cap[0,t)}\{X_q\in O\}$.
Jeżeli $\tau_O<t$, istnieje chwila $u<t$ z $X_u\in O$
(z definicji infimum). Ponieważ $O$ jest otwarty,
a trajektoria ciągła, istnieje otoczenie $u$,
w którym trajektoria nadal należy do $O$; zawiera
ono wymierną chwilę $q<t$. Implikacja odwrotna
jest oczywista. Zatem $\{\tau_O<t\}\in\calF_t$.
Dla każdego $t\ge0$ zachodzi
\[
\{\tau_O\le t\}=\bigcap_{n\ge1}\{\tau_O<t+1/n\}
\in\bigcap_{n\ge1}\calF_{t+1/n}=\calF_t,
\]
co dowodzi tezy. Założenie prawostronnej ciągłości
pojawia się dokładnie w ostatnim przejściu.
\end{proof}

\begin{uwaga}[Wersja dla procesu ciągłego tylko prawie na pewno]
Jeśli ciągłość trajektorii zachodzi na mierzalnym
zbiorze pełnej miary $\Omega_0$, to zbiory opisujące
czas trafienia mogą różnić się od zdarzeń mierzalnych
przez podzbiory $\Omega\setminus\Omega_0$. Wystarczy
pracować z ciągłą modyfikacją o trajektoriach ciągłych
wszędzie (zmieniając proces poza $\Omega_0$) albo
założyć zupełność filtracji; wówczas powyższe
argumenty zachowują ważność po uwzględnieniu
zbiorów zerowych.
\end{uwaga}

\section{Informacja dostępna w losowej chwili}
\begin{definicja}[$\sigma$-ciało w czasie zatrzymania]
Niech $\tau$ będzie czasem zatrzymania dla filtracji
$\mathbb F=(\calF_t)_{t\ge0}$ na $(\Omega,\calF,\Prob)$.
Definiujemy
\[
\calF_\tau:=\big\{A\in\calF:\ 
A\cap\{\tau\le t\}\in\calF_t\quad
\text{dla każdego }t\ge0\big\}.
\]
Interpretujemy $\calF_\tau$ jako informację dostępną
w chwili losowej $\tau$.
\end{definicja}

\begin{twierdzenie}[Podstawowe własności $\calF_\tau$]
Niech $\mathbb F$ będzie filtracją i $\tau$ jej czasem
zatrzymania. Wtedy $\calF_\tau$ jest $\sigma$-ciałem.
Jeśli $\sigma$ jest innym czasem zatrzymania i
$\sigma\le\tau$ punktowo, to
$\calF_\sigma\subseteq\calF_\tau$.
Dla każdego $A\in\calF$ zachodzi również równoważność
\[
A\in\calF_\tau
\quad\Longleftrightarrow\quad
A\cap\{\tau=t\}\in\calF_t
\ \text{dla każdego }t\in\N_0,
\]
\emph{jeżeli} $\tau$ przyjmuje wartości w $\N_0\cup\{\infty\}$.
Nie jest wtedy potrzebny osobny warunek na zdarzeniu
$\{\tau=\infty\}$, ponieważ definicja zawiera tylko
przecięcia ze zdarzeniami $\{\tau\le t\}$ dla skończonych $t$.
\end{twierdzenie}
\begin{proof}
Własność $\Omega\in\calF_\tau$ wynika z
$\{\tau\le t\}\in\calF_t$.
Jeśli $A\in\calF_\tau$, to
$A^c\cap\{\tau\le t\}
=\{\tau\le t\}\setminus(A\cap\{\tau\le t\})\in\calF_t$.
Dla ciągu $(A_n)$ elementów $\calF_\tau$ mamy
$(\bigcup_nA_n)\cap\{\tau\le t\}
=\bigcup_n(A_n\cap\{\tau\le t\})\in\calF_t$.
Tak otrzymujemy $\sigma$-ciało.

Jeśli $\sigma\le\tau$ i $A\in\calF_\sigma$, to
$A\cap\{\tau\le t\}
=(A\cap\{\sigma\le t\})\cap\{\tau\le t\}\in\calF_t$,
ponieważ $\{\tau\le t\}\subseteq\{\sigma\le t\}$.
Stąd $A\in\calF_\tau$.

Dla ostatniej równoważności, gdy $\tau$ jest
$\N_0\cup\{\infty\}$-wartościowy, zauważmy, że
$\{\tau=k\}=\{\tau\le k\}\setminus\{\tau\le k-1\}$
(z odpowiednią konwencją dla $k=0$), więc jest w $\calF_k$.
Jeśli $A\in\calF_\tau$, to
$A\cap\{\tau=k\}
=(A\cap\{\tau\le k\})\cap\{\tau=k\}\in\calF_k$.
W drugą stronę, jeśli każde $A\cap\{\tau=k\}$
jest w $\calF_k$, to dla dowolnego $t\ge0$
\[
A\cap\{\tau\le t\}
=\bigcup_{k=0}^{\lfloor t\rfloor}
(A\cap\{\tau=k\})\in\calF_t.
\]
Skończona suma nie obejmuje $\{\tau=\infty\}$,
ponieważ ta część nie należy do $\{\tau\le t\}$
dla skończonego $t$. Nie trzeba narzucać osobnego
warunku na $A\cap\{\tau=\infty\}$.
\end{proof}

\begin{uwaga}[Kiedy $\tau=\infty$ z dodatnim prawdopodobieństwem]
Definicja $\calF_\tau$ nie wymaga $\tau<\infty$
prawie na pewno. Jeżeli jednak chcemy badać zmienną
$X_\tau$, musimy osobno określić, co znaczy $X_\infty$,
albo ograniczyć się do zdarzenia $\{\tau<\infty\}$.
W literaturze często przyjmuje się $\calF=\calF_\infty$,
gdzie $\calF_\infty=\sigma(\bigcup_t\calF_t)$
(po odpowiednim uzupełnieniu). Bez takiego ustalenia
nie należy utożsamiać $\calF_\infty$ z całym $\calF$.
\end{uwaga}

\begin{twierdzenie}[Mierzalność obserwacji w chwili zatrzymania]
\label{thm:sampled}
Niech $E$ będzie przestrzenią polską,
$X=(X_t)_{t\ge0}$ progresywnie mierzalnym procesem
$E$-wartościowym, a $\tau$ czasem zatrzymania
względem $\mathbb F$. Załóżmy $\tau<\infty$
prawie na pewno i że $\calF_0$ zawiera wszystkie
zbiory zerowe (aby można było dowolnie określić
$X_\tau$ na $\{\tau=\infty\}$).
Wtedy zmienna $X_\tau(\omega):=X_{\tau(\omega)}(\omega)$
na $\{\tau<\infty\}$, dowolnie uzupełniona w
$\{\tau=\infty\}$, jest $\calF_\tau$-mierzalna.
\end{twierdzenie}
\begin{proof}
Ustalmy $t\ge0$. Z własności czasu zatrzymania
zmienna $\tau\wedge t$ jest $\calF_t$-mierzalna.
Progresywna mierzalność $X$ oznacza, że odwzorowanie
$(s,\omega)\mapsto X_s(\omega)$ na $[0,t]\times\Omega$
jest mierzalne względem
$\calB([0,t])\otimes\calF_t$.
Odwzorowanie
$\omega\mapsto(\tau(\omega)\wedge t,\omega)$
z $\Omega$ do $[0,t]\times\Omega$ jest
$\calF_t$-mierzalne, stąd
$X_{\tau\wedge t}$ jest $\calF_t$-mierzalna.
Na zdarzeniu $\{\tau\le t\}$ mamy
$X_\tau=X_{\tau\wedge t}$, a więc dla każdego
$B\in\calB(E)$ zachodzi
\[
\{X_\tau\in B\}\cap\{\tau\le t\}
=\{X_{\tau\wedge t}\in B\}\cap\{\tau\le t\}\in\calF_t.
\]
Z definicji oznacza to
$\{X_\tau\in B\}\in\calF_\tau$.
Dla $\tau=\infty$ na zbiorze zerowym uzupełnienie
jest mierzalne przy założonej zupełności filtracji.
\end{proof}

\section{Zatrzymywanie procesów i próbkowanie opcjonalne}
\begin{definicja}[Proces zatrzymany]
Niech $X=(X_t)_{t\ge0}$ będzie procesem $E$-wartościowym,
a $\tau$ czasem zatrzymania. \emph{Procesem zatrzymanym}
nazywamy $X^\tau=(X_{t\wedge\tau})_{t\ge0}$.
W szczególności dla $\tau=\infty$ mamy
$X_{t\wedge\tau}=X_t$ dla wszystkich skończonych $t$.
\end{definicja}

\begin{stwierdzenie}
Niech $X$ będzie progresywnie mierzalny i adaptowany do
filtracji $\mathbb F$, a $\tau$ czasem zatrzymania.
Wtedy proces $X^\tau$ jest adaptowany. Jeśli każda
trajektoria $X$ jest ciągła (odpowiednio prawostronnie
ciągła), to taka sama regularność zachodzi dla $X^\tau$.
\end{stwierdzenie}
\begin{proof}
Dla ustalonego $t$ zmienna $X_{\tau\wedge t}$
jest $\calF_t$-mierzalna na mocy argumentu z dowodu
twierdzenia~\ref{thm:sampled}, ponieważ $\tau\wedge t$
jest $\calF_t$-mierzalna, a proces jest progresywny.
To jest adaptacja. Dla ustalonego $\omega$ funkcja
$t\mapsto t\wedge\tau(\omega)$ jest ciągła.
Złożenie jej z ciągłą trajektorią $u\mapsto X_u(\omega)$
jest ciągłe. Jeśli ta ostatnia jest tylko prawostronnie
ciągła, to przy $t_n\downarrow t$ również
$t_n\wedge\tau\downarrow t\wedge\tau$, więc
zachodzi prawostronna ciągłość złożenia.
\end{proof}

W tomie I udowodniliśmy twierdzenie o zatrzymywaniu
martyngałów w czasie dyskretnym. W czasie ciągłym
przyda się jego odpowiednik dla czasów ograniczonych.
Przytoczymy go tu z dowodem, aby od początku pilnować
różnicy pomiędzy \emph{ograniczeniem czasu} a jedynie
jego prawie pewną skończonością.

\begin{lemat}[Jednostajna całkowalność warunkowych wartości oczekiwanych]
\label{lem:ui-cond}
Niech $Y\in L^1(\Omega,\calF,\Prob)$.
Rodzina $\{\E[Y\mid\calG]:\calG\subseteq\calF
\text{ jest pod-$\sigma$-ciałem}\}$ jest jednostajnie całkowalna.
\end{lemat}
\begin{proof}
Ustalmy $K>0$ i pod-$\sigma$-ciało $\calG$.
Rozpiszmy $Y=Y^{\le}+Y^{>}$, gdzie
$Y^{\le}=Y\ind{\{|Y|\le K/2\}}$ oraz
$Y^{>}=Y\ind{\{|Y|>K/2\}}$.
Dla $Z=\E[Y\mid\calG]$ mamy
$|\E[Y^{\le}\mid\calG]|\le K/2$ oraz
$|Z|\le K/2+\E[|Y^{>}|\mid\calG]$.
Na zdarzeniu $A=\{|Z|>K\}$ otrzymujemy
$\E[|Y^{>}|\mid\calG]>K/2$ i stąd
$|Z|\ind A\le2\E[|Y^{>}|\mid\calG]\ind A$.
Po scałkowaniu i wykorzystaniu własności
warunkowej wartości oczekiwanej dostajemy
\[
\E[|Z|\ind{\{|Z|>K\}}]
\le2\E[|Y|\ind{\{|Y|>K/2\}}]\xrightarrow[K\to\infty]{}0.
\]
Oszacowanie nie zależy od $\calG$, co dowodzi
jednostajnej całkowalności całej rodziny.
\end{proof}

\begin{lemat}[Próbkowanie opcjonalne na skończonej siatce]
\label{lem:disc-optional}
Niech $(M_k)_{k=0}^N$ będzie całkowalnym martyngałem
względem filtracji $(\calH_k)_{k=0}^N$,
a $\sigma,\tau$ czasami zatrzymania o wartościach w
$\{0,1,\ldots,N\}$, spełniającymi $\sigma\le\tau$.
Wtedy $\E[M_\tau\mid\calH_\sigma]=M_\sigma$
prawie na pewno.
\end{lemat}
\begin{proof}
Niech $A\in\calH_\sigma$. Rozpiszmy różnicę
$M_\tau-M_\sigma$ wzdłuż przyrostów:
\[
M_\tau-M_\sigma
=\sum_{k=1}^N\ind{\{\sigma<k\le\tau\}}
(M_k-M_{k-1}).
\]
Dla każdego $k$ zdarzenie
$A\cap\{\sigma<k\le\tau\}$ jest $\calH_{k-1}$-mierzalne.
Rzeczywiście,
$A\cap\{\sigma\le k-1\}\in\calH_{k-1}$
z definicji $\calH_\sigma$, a
$\{\tau\ge k\}=\{\tau\le k-1\}^c\in\calH_{k-1}$.
Własność martyngałowa daje więc
\[
\E[\ind{A\cap\{\sigma<k\le\tau\}}(M_k-M_{k-1})]
=\E\!\left[\ind{A\cap\{\sigma<k\le\tau\}}
\E[M_k-M_{k-1}\mid\calH_{k-1}]\right]=0.
\]
Sumujemy po skończenie wielu $k$ i otrzymujemy
$\E[\ind A M_\tau]=\E[\ind A M_\sigma]$.
Ponieważ $M_\sigma$ jest $\calH_\sigma$-mierzalny
(jest sumą $M_k\ind{\{\sigma=k\}}$),
charakterystyka warunkowej wartości oczekiwanej
daje tezę.
\end{proof}

\begin{twierdzenie}[Opcjonalne próbkowanie w ograniczonym horyzoncie]
\label{thm:optional-cont}
Niech $(\Omega,\calF,\Prob)$ będzie przestrzenią
probabilistyczną z filtracją $(\calF_t)_{0\le t\le T}$,
gdzie $T>0$. Niech $M=(M_t)_{0\le t\le T}$
będzie całkowalnym martyngałem względem tej filtracji
o trajektoriach prawostronnie ciągłych.
Jeśli $\sigma,\tau$ są czasami zatrzymania spełniającymi
punktowo $0\le\sigma\le\tau\le T$,
to $M_\sigma,M_\tau\in L^1$ oraz
\[
\E[M_\tau\mid\calF_\sigma]=M_\sigma
\quad\text{prawie na pewno}.
\]
\end{twierdzenie}
\begin{proof}
\emph{Krok 1: aproksymacja czasów.}
Dla każdego $n$ weźmy siatkę $D_n=\{kT2^{-n}:0\le k\le2^n\}$
i niech $\tau_n$ (odpowiednio $\sigma_n$) będzie
najmniejszym punktem siatki nie mniejszym od $\tau$
(odpowiednio $\sigma$). Wtedy
$\sigma_n\le\tau_n$, $\tau_n\downarrow\tau$,
$\sigma_n\downarrow\sigma$ i oba przybliżenia są
czasami zatrzymania o wartościach w $D_n$.
Istotnie, dla $kT2^{-n}\in D_n$ zachodzi
$\{\tau_n\le kT2^{-n}\}=\{\tau\le kT2^{-n}\}\in\calF_{kT2^{-n}}$,
podobnie dla $\sigma_n$.
Lemat~\ref{lem:disc-optional} daje
$\E[M_{\tau_n}\mid\calF_{\sigma_n}]=M_{\sigma_n}$.

\emph{Krok 2: jednostajna całkowalność.}
Stosując ten sam lemat do czasu $\tau_n$ i
stałego czasu końcowego $T$, otrzymujemy
$M_{\tau_n}=\E[M_T\mid\calF_{\tau_n}]$.
Z lematu~\ref{lem:ui-cond} rodzina
$(M_{\tau_n})$ jest jednostajnie całkowalna;
tak samo $(M_{\sigma_n})$.
Prawostronna ciągłość trajektorii wraz z
$\tau_n\downarrow\tau$, $\sigma_n\downarrow\sigma$
powoduje zbieżności
$M_{\tau_n}\to M_\tau$ i
$M_{\sigma_n}\to M_\sigma$ prawie na pewno.
Twierdzenie Vitaliego (zbieżność prawie na pewno i
jednostajna całkowalność implikują zbieżność w $L^1$)
daje obie zbieżności w $L^1$. W szczególności
zmienne graniczne są całkowalne.

\emph{Krok 3: przejście do warunkowania w losowej chwili.}
Niech $A\in\calF_\sigma$. Ponieważ
$\sigma\le\sigma_n$ punktowo, z poprzedniego
twierdzenia o informacjach w czasach zatrzymania
wynika $A\in\calF_{\sigma_n}$. Z równości
warunkowych wartości oczekiwanych z kroku 1 dostajemy
$\E[\ind A M_{\tau_n}]=\E[\ind A M_{\sigma_n}]$.
Przechodząc do granicy w $L^1$, otrzymujemy
$\E[\ind A M_\tau]=\E[\ind A M_\sigma]$.
Zmienna $M_\sigma$ jest $\calF_\sigma$-mierzalna,
ponieważ $M$ jest progresywnie mierzalny jako
proces adaptowany prawostronnie ciągły oraz na mocy
twierdzenia~\ref{thm:sampled}.
Wobec tego charakterystyka warunkowej wartości
oczekiwanej daje żądane równanie.
\end{proof}

\begin{uwaga}[Przykład pułapki: czas nieograniczony]
Twierdzenie wymaga $\tau\le T$ dla ustalonego $T$.
Dla symetrycznego spaceru losowego
$S_n=\xi_1+\cdots+\xi_n$, $\E S_n=0$,
czas $\tau=\inf\{n\ge0:S_n=1\}$ jest prawie
na pewno skończony, ale $S_\tau=1$ prawie na pewno.
Zatem $\E S_\tau=1\ne0=\E S_0$.
Nie ma sprzeczności: brak ograniczenia $\tau$
uniemożliwia bez dodatkowych hipotez przejście
od $S_{\tau\wedge n}$ do $S_\tau$ w $L^1$.
To samo ostrzeżenie będzie obowiązywać dla
martyngałów w czasie ciągłym.
\end{uwaga}


\begin{twierdzenie}[Dwa martyngały związane z procesem Wienera]
Niech $W=(W_t)_{t\ge0}$ będzie standardowym procesem
Wienera o ciągłych trajektoriach,
a $\calF_t^W=\sigma(W_s:0\le s\le t)$ jego filtracją
naturalną. Wtedy procesy
$M_t=W_t$ oraz $N_t=W_t^2-t$, $t\ge0$,
są całkowalnymi martyngałami względem $\mathbb F^W$.
\end{twierdzenie}
\begin{proof}
Ustalmy $0\le s<t$ i połóżmy
$Z=W_t-W_s$. Wiemy, że $Z\sim\Normal(0,t-s)$
i jest niezależne od wektora
$(W_{u_1},\ldots,W_{u_m})$ dla dowolnych
$u_1,\ldots,u_m\le s$: wynika to z gaussowskości
wspólnego wektora i zerowych kowariancji $Z$
z każdą wcześniejszą współrzędną.
Aby przejść od skończenie wielu współrzędnych do
całego $\sigma$-ciała $\calF_s^W$, stosujemy
twierdzenie $\pi$--$\lambda$: niezależność zachodzi
na $\pi$-układzie zdarzeń cylindrycznych
$\{(W_{u_1},\ldots,W_{u_m})\in B\}$, które generują
$\calF_s^W$, a klasa zdarzeń niezależnych od $Z$
jest $\lambda$-układem. Zatem $Z$ jest niezależne
od $\calF_s^W$.

Ponieważ $\E Z=0$ i $\E Z^2=t-s$, z własności
warunkowej wartości oczekiwanej dostajemy
$\E[Z\mid\calF_s^W]=0$ oraz
$\E[Z^2\mid\calF_s^W]=t-s$.
Z $W_t=W_s+Z$ otrzymujemy
\[
\E[W_t\mid\calF_s^W]=W_s,
\quad
\E[W_t^2-t\mid\calF_s^W]
=W_s^2+2W_s\E[Z\mid\calF_s^W]
+(t-s)-t=W_s^2-s.
\]
Całkowalność wynika ze skończonych momentów
rozkładu normalnego, a adaptacja z definicji
filtracji naturalnej. To są wszystkie warunki
bycia martyngałem.
\end{proof}

\begin{zadanie}[Wyjście procesu Wienera z przedziału]
Niech $W$ będzie procesem Wienera z ciągłymi trajektoriami,
$W_0=0$, a $a>0$. Względem jego filtracji naturalnej
zdefiniujmy
$\tau:=\inf\{t\ge0:|W_t|=a\}$.
Wykażemy, że $\tau<\infty$ prawie na pewno,
obliczymy $\E\tau$ oraz prawdopodobieństwa
wyjścia z przedziału przez jego prawy i lewy koniec.

Najpierw $\tau$ jest czasem zatrzymania,
ponieważ jest czasem trafienia w domknięty zbiór
$\{-a,a\}$ przez ciągły adaptowany proces.
Dla $T>0$ zastosujmy twierdzenie o próbkowaniu
opcjonalnym do ograniczonego czasu
$\tau\wedge T$ i obu martyngałów
$W_t$ i $W_t^2-t$. Otrzymujemy
\[
\E W_{\tau\wedge T}=0,
\qquad \E W_{\tau\wedge T}^2=\E(\tau\wedge T).
\tag{8}
\]
Z ciągłości trajektorii i definicji pierwszego
wyjścia wynika $|W_{\tau\wedge T}|\le a$
punktowo. Wobec tego z (8) mamy
$0\le\E(\tau\wedge T)\le a^2$.
Gdy $T\uparrow\infty$, twierdzenie o zbieżności
monotonicznej daje $\E\tau\le a^2<\infty$.
W szczególności $\tau<\infty$ prawie na pewno:
gdyby $\Prob(\tau=\infty)>0$, całka $\E\tau$
byłaby nieskończona.

Skoro $\tau$ jest skończony i trajektorie są ciągłe,
mamy $|W_\tau|=a$ prawie na pewno.
Ponadto $W_{\tau\wedge T}\to W_\tau$ prawie na pewno
i $|W_{\tau\wedge T}|\le a$, więc przez zbieżność
zdominowaną $\E W_\tau=0$ oraz
$\E W_{\tau\wedge T}^2\to\E W_\tau^2=a^2$.
Przechodząc do granicy w (8), dostajemy
$\E\tau=a^2$.

Na koniec oznaczmy $p=\Prob(W_\tau=a)$.
Ponieważ $W_\tau\in\{-a,a\}$ prawie na pewno,
$0=\E W_\tau=ap-a(1-p)$, skąd $p=1/2$.
Zatem cząstka wychodzi przez każdy z obu końców z
prawdopodobieństwem $1/2$, a oczekiwany czas wyjścia
z przedziału $(-a,a)$ wynosi $a^2$.

Ten przykład ilustruje właściwe użycie twierdzenia
o zatrzymywaniu: \emph{najpierw} zatrzymujemy proces
w chwili $\tau\wedge T$, gdzie czas jest ograniczony,
a dopiero potem uzasadniamy przejście $T\to\infty$
za pomocą oszacowania i odpowiednich twierdzeń
o zbieżności.
\end{zadanie}

\section{Zadanie rozwiązane: informacja, zatrzymanie i uczciwa gra}
\begin{zadanie}[Dwa rzuty monetą i chwila zatrzymania]
Niech $\Omega=\{HH,HT,TH,TT\}$, gdzie litera $H$
oznacza wynik $+1$, a $T$ wynik $-1$.
Wszystkim czterem wynikom przypisujemy prawdopodobieństwo $1/4$.
Niech $\xi_1,\xi_2$ będą odpowiednimi wynikami
rzutów, $\calF_0=\{\varnothing,\Omega\}$,
$\calF_1=\sigma(\xi_1)$, $\calF_2=\sigma(\xi_1,\xi_2)$,
$S_0=0$, $S_1=\xi_1$ i $S_2=\xi_1+\xi_2$.
Zdefiniujmy $\tau=1$ wtedy i tylko wtedy, gdy
$\xi_1=+1$, oraz $\tau=2$ w przeciwnym razie.

\emph{(a) Sprawdzić, że $\tau$ jest czasem zatrzymania.}
Dla $t<1$ zdarzenie $\{\tau\le t\}$ jest puste.
Dla $1\le t<2$ jest równe $\{\xi_1=+1\}$,
a więc należy do $\calF_1\subseteq\calF_t$
(przy stałym przedłużeniu filtracji między rzutami).
Dla $t\ge2$ jest równe $\Omega$. Stąd warunek czasu
zatrzymania jest spełniony.

\emph{(b) Opisać dokładnie $\calF_\tau$.}
Wyniki $HH$ i $HT$ są nierozróżnialne w chwili
zatrzymania: po pierwszym $H$ zatrzymujemy się,
więc drugiego rzutu nie obserwujemy. Gdy pierwszy
rzut da $T$, czekamy do drugiego rzutu i rozróżniamy
$TH$ oraz $TT$. Zatem atomami $\calF_\tau$ są
$\{HH,HT\}$, $\{TH\}$ i $\{TT\}$.
Można to sprawdzić z definicji: warunek
$A\cap\{\tau\le1\}\in\calF_1$ pozwala
albo wziąć obydwa wyniki $HH,HT$, albo żaden;
na $\{\tau=2\}$ nie ma dodatkowego ograniczenia,
bo $\calF_2=2^\Omega$.

\emph{(c) Obliczyć rozkład i średnią $S_\tau$.}
Dla $HH$ i $HT$ mamy $\tau=1$ oraz $S_\tau=1$.
Dla $TH$ mamy $\tau=2$ i $S_\tau=0$,
a dla $TT$ mamy $\tau=2$ i $S_\tau=-2$. Stąd
\[
\Prob(S_\tau=1)=\tfrac12,\qquad
\Prob(S_\tau=0)=\tfrac14,\qquad
\Prob(S_\tau=-2)=\tfrac14,
\]
co daje $\E S_\tau=\tfrac12\cdot1+
\tfrac14\cdot0+\tfrac14\cdot(-2)=0$.
Jest to konkretny przykład próbkowania opcjonalnego:
$S$ jest martyngałem, a $\tau\le2$.

\emph{(d) Co będzie, jeśli wybierzemy czas na podstawie przyszłości?}
Niech $\rho=1$ dla $\xi_2=+1$ oraz $\rho=2$
dla $\xi_2=-1$. Zdarzenie $\{\rho\le1\}
=\{\xi_2=+1\}$ nie należy do $\calF_1$,
więc $\rho$ nie jest czasem zatrzymania.
Mamy $S_\rho=\xi_1+\xi_2\ind{\{\xi_2=-1\}}$,
zatem $\E S_\rho=0-\tfrac12=-\tfrac12$.
Nie narusza to twierdzenia o próbkowaniu opcjonalnym:
strategia używa przyszłego wyniku, czyli jej czas nie
spełnia kluczowego założenia.
\end{zadanie}

\section{Co będzie potrzebne w rachunku Itô?}
Na zakończenie zbierzmy logiczne zależności.
Aby konstruować całkę stochastyczną względem procesu Wienera,
będziemy musieli odróżniać następujące własności:
\begin{itemize}
\item proces jest \emph{adaptowany}, a więc stan w chwili $t$
jest znany przy informacji $\calF_t$;
\item proces jest \emph{progresywnie mierzalny}, co umożliwia
mierzalne podstawianie losowych chwil i zwykłe całkowanie
funkcji $(t,\omega)$;
\item proces jest \emph{przewidywalny}, co pozwala zbudować
całkę z wartości współczynnika znanych przed
przyrostem $W_{t_{k+1}}-W_{t_k}$;
\item proces jest \emph{zatrzymany} w czasie $\tau$,
co pozwala lokalizować oszacowania i badać zjawiska
do losowej chwili, bez wiedzy o przyszłości.
\end{itemize}
W dalszych rozdziałach omówimy proces Wienera
wraz z jego naturalną filtracją, procesy Markowa
w czasie ciągłym, martyngały lokalne, miarę Wienera
na przestrzeni trajektorii i całkę Itô.


% ================================================================
% TOM II. ROZDZIAŁ 4
% ================================================================
